\section{Experimental Design}

\subsection{Questions tested}

The experiments are organized around five questions.

First, if the input spectrum is clean, can the spectral-geometry calibrator
distinguish linear monofractal, boundary MRW, curved MRW, and unstable spectra?
Second, on finite raw samples, do monofractal and non-MRW controls create false
MRW-compatible evidence? Third, can deterministic estimators recover the MRW
curvature coordinate \(\lambda^2\) under the tested sample lengths and scale
ranges? Fourth, how much zeta-space noise is needed to destroy the analytic
separation between linear and curved spectra? Fifth, can the framework be run
on real complex-system data as a cautious diagnostic sanity check?

This ordering is deliberate. The paper does not start from the assumption that
MRW parameters are recoverable. It first checks the analytic spectral
geometry and then tests where this geometry breaks down under finite-sample
estimation.

\subsection{Synthetic process families}

The controlled raw-sample experiments use the following process families:
MRW, low-\(\lambda^2\) MRW, shuffled MRW, fractional Gaussian noise, iid
Gaussian, iid Student-\(t\), GARCH(1,1), and regime-switching Gaussian. The
MRW and fGn families provide interpretable scaling targets. Gaussian, Student
\(t\), GARCH, and regime-switching controls are included because heavy tails,
volatility clustering, and nonstationary variance can produce apparent
multiscale behavior without implying an MRW mechanism
\cite{bollerslev1986garch,hamilton1989regime}.

The default diagnostic \(q\)-grid is
\(\{0.5,1.0,1.5,2.0,2.5,3.0\}\). Process-family diagnostics use lengths
\(T=256,512,1024,2048\). The deterministic identifiability supplement reports
short-window grids including \(T=512,1024,2048\), with scale sets based on
powers of two. All such claims are restricted to the tested grids.

\subsection{Analytic spectrum-space experiments}

The analytic spectrum-space dataset does not generate raw time series. It
directly generates \(\zeta(q)\) curves from known families: linear
monofractal, boundary MRW, curved MRW, noisy monofractal, noisy MRW,
heavy-tail distorted, regime-apparent, and ambiguous mild-curvature spectra.
This experiment tests only the interpretation layer. It answers whether
linear, boundary, and curved geometries are separable once the spectrum is
given.

To check that the analytic result is not seed-specific, the evaluation was
repeated for seeds 2024, 2025, and 2026. The resulting mean and standard
deviation tables are included in the paper assets.
Training configurations, output tables, and figure data are kept in the
repository so that the reported diagnostics can be regenerated without relying
on undocumented manual processing.

\subsection{Deterministic and classical baselines}

The finite-sample identifiability study compares several deterministic
structure-function variants: ordinary least squares, trimmed regression,
bootstrap over scales, and smoothed \(\zeta(q)\). In addition, classical
aggregate-increment structure functions, first- and second-order MFDFA, and
dependency-free Haar wavelet-leader / WTMM-style modulus-maxima baselines were
run on the same quick finite-sample grid. These baselines are not used to
claim that CMIN-SR outperforms all classical methods. They are used to test
whether representative classical estimators remove the short-window
\(\lambda^2\) recovery bottleneck in the tested settings
\cite{muzy1993multifractal,kantelhardt2002mfdfa,jaffard2007wavelet,wendt2007bootstrap}.

The Haar wavelet-leader and WTMM-style rows should be interpreted as compact
classical controls, not as full production wavelet-leader or continuous-WTMM
packages. They are included to test whether lightweight classical alternatives
remove the observed bottleneck, not to close the broader benchmarking question.

\subsection{Zeta-noise bridge and failure attribution}

The zeta-noise bridge adds controlled perturbations to analytic spectra before
feeding them to the spectral-geometry calibrator. Smooth, high-\(q\), jagged,
and high-\(q\)-biased perturbations are included. This experiment measures how
the calibrated MRW-vs-linear separation margin decreases as the input spectrum
becomes less reliable.

The failure-attribution experiment compares three levels for the same process
families: analytic target spectra, deterministic estimated spectra, and neural
empirical spectra. This directly tests whether the bottleneck lies in the
spectral interpretation layer or in the upstream empirical spectrum.

\subsection{Real-world sanity check}

As an exploratory real-data check, the diagnostic proxy pipeline is applied to
local Fama-French factor returns. The four selected factor series are Mkt-RF,
SMB, HML, and Mom. Rolling windows of length 512 are compared against shuffled
surrogates. The goal is modest: to verify that the diagnostic shows
original-vs-shuffled differences consistent with temporal multiscale
dependence. These results
are not used to claim that financial factor returns are generated by an MRW
process.

\subsection{Evaluation metrics}

The main reported quantities are diagnostic scores
\((p_{\mathrm{scaling}},p_{\mathrm{curved}},p_{\mathrm{mono}},
p_{\mathrm{MRW}},p_{\mathrm{boundary}})\), projection residuals,
MRW-vs-monofractal gain, \(\lambda^2_{\mathrm{proj}}\), zeta error, and
finite-sample recovery summaries such as \(\lambda^2\) MAE, RMSE, correlation,
Spearman correlation, and high-\(\lambda^2\) detection rate. For real-data
sanity checks, the reported quantities are original-minus-shuffled gaps in
proxy \(\lambda^2\), MRW diagnostic score, spectrum width, and log-volatility
covariance slope.
